% Part: proof-theory
% Chapter: sequent-calculus
% Section: translations

\documentclass[../../../include/open-logic-section]{subfiles}

\begin{document}

\olfileid{pt}{seq}{tr}

\olsection{Ngowahi bukti antarane \Log{G1c} lan \Log{G3c}}

Kita wis nyebutake yen \Log{G1c} lan \Log{G3c} kalorone sahih lan
jangkep kanggo logika klasik, yaiku saben kalkulus bisa mbuktekake kabeh
lan mung sekuen kang bener ing saben !!{structure}. Mligine, kalorone
mbuktekake sekuen kang padha. Saiki kita bisa mbuktekake kasunyatan iki
kanthi cara teori bukti, tanpa liwat teorema kasahihan lan kajangkepan.
Bukti kaya mangkono menehi \emph{pangowahan} saka !!{proof}s ing
\Log{G1c} dadi !!{proof}s ing \Log{G3c}, lan uga kosok baline.

\begin{prop}
Yen $\Log{G1c} \Proves S$, mula $\Log{G3c} \Proves S$.
\end{prop}

\begin{proof}
  Kita nuduhake yen $\pi$ minangka !!a{proof} ing \Log{G1c}, mula ana
  !!a{proof}~$\pi'$ kanggo~$S$ ing \Log{G3c}.
  Kita nindakake iki kanthi induksi marang $n = \pheight{\pi}$.

Yen $n = 0$, mula $\pi$ minangka aksioma. Aksioma
$\lfalse \Sequent \quad$ uga minangka aksioma ing~\Log{G3c}
(pilih konteks $\Gamma, \Delta$ kosong). Ing \Log{G3c}, aksioma
$!A \Sequent !A$ ora diidini kanggo sembarang~$!A$;
formula kasebut kudu atomik. Nanging ing \olref[ex]{prop:G3c-axioms}
kita wis mbuktekake yen kabeh sekuen kaya mangkono !!{provable}
ing~\Log{G3c}.

Yen $n>0$, $\pi$ dipungkasi nganggo aturan~$R$.
Hipotesis induktif njamin ana !!{proof}s kanggo premis aturan kasebut
ing~\Log{G3c}, yaiku !!{proof}s bagean langsung nduweni pangowahan
ing~\Log{G3c}. Yen aturan~$R$ uga minangka aturan ing \Log{G3c},
kita ngetrapake aturan kasebut marang !!{proof}s premis kang wis diowahi.
Kanggo aturan kang dudu aturan ing~\Log{G3c}, kita wis nuduhake
yen aturan kasebut paling ora admisibel. Mula, yen aturan~$R$
yaiku $\LeftR{\land}$ utawa $\RightR{\lor}$,
gunakna \olref[adm]{prop:lor-G3-adm}. Yen minangka aturan pelemahan,
gunakna \olref[adm]{prop:weak-G3c-adm}. Yen minangka aturan kontraksi,
gunakna \olref[inv]{prop:G3c-cont-adm}.
Kanggo \LeftR{\lforall} lan \RightR{\lexists} ing \Log{G1c},
gunakna \LeftR{\Weakening} utawa \RightR{\Weakening} kang admisibel
dhisik kanggo nambah formula utama kanthi kuantor kang kudu tetep ana,
banjur trapna aturan \Log{G3c} kang cocog.
\end{proof}

\begin{prop}
Yen $\Log{G3c} \Proves S$, mula $\Log{G1c} \Proves S$.
\end{prop}

\begin{proof}
Kita maneh nggunakake induksi marang !!{height} saka !!a{proof}~$\pi$
kanggo $S$. Saben aksioma \Log{G3c} bisa !!{prove} ing~\Log{G1c}
saka sawijining aksioma kanthi ngetrapake aturan \LeftR{\Weakening}
lan \RightR{\Weakening} kaping pirang-pirang:
\[
  \Axiom$!C \fCenter !C$
  \RightLabel{\LeftR{\Weakening}}
  \Deduce$!C, \Gamma \fCenter !C$
  \RightLabel{\RightR{\Weakening}}
  \Deduce$!C, \Gamma \fCenter \Delta, !C$
  \DisplayProof
\qquad
  \Axiom$\lfalse \fCenter$
  \RightLabel{\LeftR{\Weakening}}
  \Deduce$\lfalse, \Gamma \fCenter$
  \RightLabel{\RightR{\Weakening}}
  \Deduce$\lfalse, \Gamma \fCenter \Delta$
  \DisplayProof
  \]
Yen $\pi$ dipungkasi nganggo aturan~$R$, hipotesis induktif menehi
!!{proof}s ing \Log{G1c} kanggo premis kasebut, lan kita bisa ngetrapake
aturan~$R$ kang padha marang !!{proof}s kasebut.
Pangecualiane yaiku aturan \LeftR{\land}, \RightR{\lor},
\LeftR{\lforall}, lan \RightR{\lexists}, kang beda antarane
\Log{G1c} lan~\Log{G3c}. Varian \Log{G3c} bisa diturunake
ing \Log{G1c} kaya mangkene:
\[
  \Axiom$!A, !B, \Gamma \fCenter \Delta$
  \RightLabel{\LeftR{\land}}
  \UnaryInf$!A \land !B, !B, \Gamma \fCenter \Delta$
  \RightLabel{\LeftR{\land}}
  \UnaryInf$!A \land !B, !A \land !B, \Gamma \fCenter \Delta$
  \RightLabel{\LeftR{\Contraction}}
  \UnaryInf$!A \land !B, \Gamma \fCenter \Delta$
    \DisplayProof
\qquad
  \Axiom$\Gamma \fCenter \Delta, !A, !B$
  \RightLabel{\RightR{\lor}}
  \UnaryInf$\Gamma \fCenter \Delta, !A \lor !B, !B$
  \RightLabel{\RightR{\lor}}
  \UnaryInf$\Gamma \fCenter \Delta, !A \lor !B, !A \lor !B$
  \RightLabel{\RightR{\Contraction}}
  \UnaryInf$\Gamma \fCenter \Delta, !A \lor !B$
    \DisplayProof
  \]
Formula utama kang tetep ana ing \LeftR{\lforall} lan \RightR{\lexists}
ditangani ing \Log{G1c} kanthi aturan kang cocog banjur kontraksi:
\[
  \Axiom$!A(t), \lforall[x][!A(x)], \Gamma \fCenter \Delta$
  \RightLabel{\LeftR{\lforall}}
  \UnaryInf$\lforall[x][!A(x)], \lforall[x][!A(x)], \Gamma \fCenter \Delta$
  \RightLabel{\LeftR{\Contraction}}
  \UnaryInf$\lforall[x][!A(x)], \Gamma \fCenter \Delta$
  \DisplayProof
  \qquad
  \Axiom$\Gamma \fCenter \Delta, \lexists[x][!A(x)], !A(t)$
  \RightLabel{\RightR{\lexists}}
  \UnaryInf$\Gamma \fCenter \Delta, \lexists[x][!A(x)], \lexists[x][!A(x)]$
  \RightLabel{\RightR{\Contraction}}
  \UnaryInf$\Gamma \fCenter \Delta, \lexists[x][!A(x)]$
  \DisplayProof
  \]
\end{proof}

\end{document}
